\begin{textsample}{POS Dim 2 – vem\_ed\_31 – Score 11.00 – vem\_ed\_31/t163.txt}  \label{ex:f2_pos_014}
Preparando alunos para profissões que não existem “Prepping COIL students for jobs not yet created”.
Os projetos COIL (Collaborative Online International Learning, ou Aprendizagem Colaborativa Internacional On-line) \textbf{contribuindo} para preparar estudantes para empregos que ainda não foram criados.
Esse foi o tema do SUNY COIL Community Webinar \textbf{conduzido} em 17 de setembro de 2025 por Osvaldo Succi Junior, coordenador da equipe de Apoio à Internacionalização do Ensino Superior da Depes / CGESG / CPS, e Luciane Stallivieri, da Universidade Federal de Santa Catarina.
Luciane Stallivieri propõe a noção de Internacionalização Responsável, que busca equilíbrio (linguístico, geográfico, econômico) na escolha de parceiros, compartilha resultados com a sociedade (accountability), oferece inclusão e garante \textbf{sustentabilidade} aos processos, com comprometimento institucional (compliance).
Para encarar o “VUCA world” (VUCA é um acrônimo em inglês para volatilidade, incerteza, complexidade e ambiguidade) é \textbf{preciso} ter visão / foco no propósito, compreensão para refinar / adaptar planos, clareza na criação desses planos e agilidade / flexibilidade e colaboração para alcançar os objetivos.
Considerando as macrotendências das transformações no mundo dos negócios, as incertezas, as mudanças aceleradas e os empregos a serem criados (e destruídos), Osvaldo Succi defende que a educação superior \textbf{deve} evoluir para preparar os estudantes com habilidades que transcendam as \textbf{descrições} específicas de \textbf{cargos}, promovendo adaptabilidade, pensamento crítico e envolvimento ético global.
Afinal, as principais competências centrais no mundo atual são: pensamento analítico; resiliência, flexibilidade e agilidade; liderança; pensamento \textbf{criativo}; \textbf{motivação} e autoconsciência.
Após um exercício \textbf{criativo} com profissões que ainda não existem, como “corretor-guardião de dados pessoais” ou “terapeuta de robôs”, Osvaldo Succi listou algumas competências desejáveis a serem incorporadas no \textbf{design} de projetos COIL (Collaborative Online International Learning, ou Aprendizagem Colaborativa Internacional On-line): Sensibilidade lógica humano-computador Resumo confiável de conteúdo \textbf{Design} ético e auditoria de IA Avaliação de riscos Desenvolvimento de cenários plausíveis Escalabilidade contextual Aprendizagem e desaprendizagem adaptativas Pensamento sistêmico Gestão da privacidade

% matched lemmas: cargo, conduzir, contribuir, criativo, descrição, design, dever, motivação, preciso, sustentabilidade
\end{textsample}
